\todo{Lecture 16}

Autant des lignes du champ $\boldsymbol{E}$ entrent et sortent de $S$, donc le flux est nul. Avec les deux surfaces $S_{1}$ et $S_{2}$, on trouve 
\begin{align*}
\iint _{S_{1}+S_{2}}\boldsymbol{E}~d\boldsymbol{S} & \overset{\mathrm{cas}~2}= \frac{q}{\varepsilon_{0}} \\
  & = \underbrace{ \iint _{S_{1}}\boldsymbol{E}~d\boldsymbol{S} }_{\overset{d\to 0}\longrightarrow \frac{q}{\varepsilon_{0}}}+\underbrace{ \iint _{S_{2}}\boldsymbol{E}~d\boldsymbol{S} }_{ \overset{d\to 0}\longrightarrow \iint _{S}\boldsymbol{E}~d\boldsymbol{S} } \\
\Rightarrow \iint _{S}\boldsymbol{E}~d\boldsymbol{S} & =0
\end{align*}
et la loi de Gauss est satisfaite dans ce cas.

Le cas général pour une densité de charge arbitraire (qui est une somme des charges ponctuelles), suit du principe de superposition.

\section{Le potentiel électrostatique}
\begin{definition}
En électrostatique, on peut définir une fonction scalaire $\phi(\boldsymbol{r})$ telle que $\boldsymbol{E}=-\boldsymbol{\nabla}\phi$. On appelle $\phi$ le potentiel électrostatique.
\end{definition}

Pour une charge ponctuelle $q_{1}$ a $\boldsymbol{r}_{1}$ on a 
$$
\phi(\boldsymbol{r})=\frac{1}{4\pi\varepsilon _{0}} \frac{q_{1}}{\lVert \boldsymbol{r}-\boldsymbol{r}_{1} \rVert }+C
$$
\begin{proof}
On a
\begin{align*}
-\boldsymbol{\nabla}\phi(\boldsymbol{r}) & =\begin{pmatrix}
-\partial _{x}\left( \frac{1}{4\pi\varepsilon_{0}} \frac{q_{1}}{\sqrt{ (x-x_{1})^{2}+(y-y_{1})^{2}+(z-z_{1})^{2} }} \right) \\
-\partial _{y}\left( \frac{1}{4\pi\varepsilon_{0}} \frac{q_{1}}{\sqrt{ (x-x_{1})^{2}+(y-y_{1})^{2}+(z-z_{1})^{2} }} \right) \\
-\partial _{z}\left( \frac{1}{4\pi\varepsilon_{0}} \frac{q_{1}}{\sqrt{ (x-x_{1})^{2}+(y-y_{1})^{2}+(z-z_{1})^{2} }} \right)
\end{pmatrix} \\
 & =\begin{pmatrix}
-\frac{1}{4\pi\varepsilon_{0}} q_{1}\left( -\frac{1}{2} \right) \frac{1}{((x-x_{1})^{2}+(y-y_{1})^{2}+(z-z_{1})^{2} )^{3/2}}\cdot2(x-x_{1}) \\
-\frac{1}{4\pi\varepsilon_{0}} q_{1}\left( -\frac{1}{2} \right) \frac{1}{((x-x_{1})^{2}+(y-y_{1})^{2}+(z-z_{1})^{2} )^{3/2}}\cdot 2(y-y_{1}) \\
-\frac{1}{4\pi\varepsilon_{0}} q_{1}\left( -\frac{1}{2} \right) \frac{1}{((x-x_{1})^{2}+(y-y_{1})^{2}+(z-z_{1})^{2} )^{3/2}}\cdot 2 (z-z_{1})
\end{pmatrix} \\
 & =\frac{q_{1}}{4\pi\varepsilon_{0}} \frac{1}{\lVert \boldsymbol{r}-\boldsymbol{r}_{1} \rVert ^{3}} \begin{pmatrix}
x-x_{1} \\
y-y_{1}  \\
z-z_{1}
\end{pmatrix}  \\
 & =\frac{q_{1}}{4\pi \varepsilon_{0}} \frac{\boldsymbol{r}-\boldsymbol{r}_{1}}{\lVert \boldsymbol{r}-\boldsymbol{r}_{1} \rVert^{3} }=\boldsymbol{E}(\boldsymbol{r})
\end{align*}
\end{proof}

Pour une distribution de charge arbitraire, on a
$$
\phi(\boldsymbol{r})=\frac{1}{4\pi\varepsilon_{0}}\iiint _{V} \frac{\rho _{\mathrm{el}}(\boldsymbol{r}')}{\lVert \boldsymbol{r}-\boldsymbol{r}' \rVert }~dV'+C
$$
\begin{proof}
On a 
\begin{align*}
-\boldsymbol{\nabla}\phi(\boldsymbol{r}) & =\frac{1}{4\pi\varepsilon_{0}}\iiint _{V}\rho _{\mathrm{el}} (\boldsymbol{r}')(-\boldsymbol{\nabla}) \frac{1}{\lVert \boldsymbol{r}-\boldsymbol{r}' \rVert  }~dV'  \\
 & =\frac{1}{4\pi \varepsilon_{0} }\iiint _{V} \frac{\rho _{\mathrm{el}}(\boldsymbol{r}')(\boldsymbol{r}-\boldsymbol{r}')}{\lVert \boldsymbol{r}-\boldsymbol{r}' \rVert ^{3}}~dV'=\boldsymbol{E}(\boldsymbol{r}).
\end{align*}
\end{proof}

Pour une densité de charge de surface :
$$
\phi (\boldsymbol{r})=\frac{1}{4\pi\varepsilon_{0}}\iint _{S} \frac{\sigma _{\mathrm{el}}(\boldsymbol{r}')}{\lVert \boldsymbol{r}-\boldsymbol{r}' \rVert }~dS'+C
$$

\subsection{Quelques propriétés de \texorpdfstring{$\phi$ et $\boldsymbol{E}$}{le potentiel électrostatique et le champ} en électrostatique}
\begin{enumerate}
\item Les unités de $\phi$ sont les volts $\mathrm{V}=\frac{\mathrm{kgm}^{2}}{\mathrm{s}^{2}}\cdot \frac{1}{\mathrm{C}}$. On l'interprète comme l'énergie par charge.
\item $\phi$ seulement définie à une constante près $\boldsymbol{E}=-\boldsymbol{\nabla}\phi$, donc $\boldsymbol{E}=-\boldsymbol{\nabla}(\phi +C)$.
\item Le potentiel électrostatique dû à différentes charges est additive (principe de superposition). Soient $\boldsymbol{E}_{1}=-\boldsymbol{\nabla}\phi_{1}$ et $\boldsymbol{E}_{2}=-\boldsymbol{\nabla}\phi_{2}$, alors $\boldsymbol{E}=-\boldsymbol{\nabla}(\phi_{1} +\phi_{2})$.
\item La rotation d'un gradient est nulle : 
$$
\boldsymbol{\nabla}\times \boldsymbol{\nabla}\phi=\begin{pmatrix}
\partial _{x} \\
\partial _{y} \\
\partial _{z}
\end{pmatrix}\times \begin{pmatrix}
\partial _{x}\phi \\
\partial _{y}\phi \\
\partial _{z}\phi 
\end{pmatrix}=0
$$
donc $\boldsymbol{\nabla}\times \boldsymbol{E}=0$.
\item L'inverse de 4 est vrai aussi : Soit un champ vectoriel $\boldsymbol{A}(\boldsymbol{r})$ avec $\boldsymbol{\nabla}\times \boldsymbol{A}=0$ pour tout $\boldsymbol{r}$. Alors, il existe $f(\boldsymbol{r})$ telle que $\boldsymbol{A}=-\boldsymbol{\nabla}f$.
\item Le travail pour déplacer une charge $q$ dans un champ $\boldsymbol{E}$ de $A$ a $B$ ne dépend pas du chemin choisi et est égal à $q(\phi(B)-\phi(A))$.
\begin{proof}
Soit $\gamma$ un chemin paramétré par $\boldsymbol{\ell}(s)$ avec $s \in[0,1]$. Alors 
\begin{align*}
W_{\gamma} & =-\int _{_{\gamma}}\boldsymbol{F}~d\boldsymbol{\ell} \\
 & =-\int _{\gamma}q\boldsymbol{E}~d\boldsymbol{\ell} \\
 & =\int _{\gamma}q\boldsymbol{\nabla}\phi~d\boldsymbol{\ell} \\
 & =\int_{0}^{1} q\boldsymbol{\nabla}\phi(\boldsymbol{\ell}(s)) \, \frac{d\boldsymbol{\ell}}{ds}ds \\
 &  =\int_{0}^{1} q~ \frac{d}{ds}(\phi(\boldsymbol{\ell}(s))) \, ds \\
 & = q\phi(\boldsymbol{\ell}(s))|_{0}^{1} \\
 & =q(\phi(B)-\phi(A))
\end{align*}
donc $W_{\gamma}$ est independant du chemin choisi entre $A$ et $B$.

Si $A=B$, il suit que le travail pour déplacer une charge $q$ le long d'une courbe fermée est nul. Ceci implique que $q\boldsymbol{E}$ est une force conservative.  
\end{proof}
\todo{Lecture 17}
\item En électrostatique, pour une charge ponctuelle $(q,m)$ et gravite $g\boldsymbol{\hat{e}}_{z}$. L'énergie $E_{\mathrm{tot}}=\frac{1}{2}m\lVert \boldsymbol{v} \rVert^{2}+q\phi(\boldsymbol{x})+mgz$ est conservée le long de la trajectoire $\boldsymbol{x}(t)$ de la particule.
\begin{proof}
\begin{align*}
\frac{d}{dt}E_{\mathrm{ t ot}} & =\frac{1}{2}m2\boldsymbol{v}\dot{\boldsymbol{v}}+q\boldsymbol{\nabla}\phi   \frac{d\boldsymbol{x}}{dt}+mg  \frac{dz}{dt} \\
 & =m \dot{\boldsymbol{v}}\boldsymbol{v} +q\boldsymbol{\nabla}\phi \boldsymbol{v}+mg\boldsymbol{\hat{e}}_{z}\boldsymbol{v} \\
 & =\boldsymbol{F}\boldsymbol{v}-(q\boldsymbol{E}-mg\boldsymbol{\hat{e}}_{z})\boldsymbol{v}=0
\end{align*}
Donc $\phi(B)-\phi(A)=\int _{A\to B}\boldsymbol{E}~d\boldsymbol{\ell}$ est le changement de l'énergie potentielle électrostatique par unite de charge entre $A$ et $B$.  $\phi(B)-\phi(A)$ est aussi appelle la tension entre $A$ et $B$.
\end{proof}
\item On définit des surfaces équipotentielles par $\phi(\boldsymbol{r})$ constante. Comme $\boldsymbol{E}=-\boldsymbol{\nabla}\phi$, les lignes du champ électrique sont partout perpendiculaires aux surfaces équipotentielles.
\begin{proof}
Par une expansion de Taylor 
\begin{align*}
\phi(\boldsymbol{r}+\boldsymbol{h}) & =\phi(\boldsymbol{r})+\boldsymbol{\nabla}\phi(\boldsymbol{r})\boldsymbol{h}+O(\lVert \boldsymbol{h} \rVert ^{2}) \\
 & =\phi(\boldsymbol{r})-\boldsymbol{E}(\boldsymbol{r})\boldsymbol{h}+O(\lVert \boldsymbol{h} \rVert ^{2})  \\
\end{align*}
Si $\boldsymbol{h}$ est tangentielle a la surface équipotentielle au point $\boldsymbol{r}$ : 
\begin{align*}
\phi(\boldsymbol{h}+\boldsymbol{r}) & =\phi(\boldsymbol{r})+O(\lVert \boldsymbol{h} \rVert ^{2}) \\
 & =\phi(\boldsymbol{r})-\boldsymbol{E}(\boldsymbol{r})\boldsymbol{h}+O(\lVert \boldsymbol{h} \rVert ^{2}) \\
 \Rightarrow \boldsymbol{E}(\boldsymbol{r})\boldsymbol{h} & =0
\end{align*}
car $\phi$ est constante sur la surface équipotentielle. Donc $\boldsymbol{E}$ pointe dans la direction perpendiculaire aux surfaces équipotentielles. 
\end{proof}
\end{enumerate}

